\documentclass[11pt, a4paper]{article}

% ──────────────────────────────────────────────
% Packages
% ──────────────────────────────────────────────
\usepackage{fontspec}
\usepackage[margin=1in, top=1in, bottom=1in]{geometry}
\usepackage{graphicx}
\usepackage{float}
\usepackage{xcolor}
\usepackage{hyperref}
\usepackage{listings}
\usepackage{booktabs}
\usepackage{longtable}
\usepackage{tabularx}
\usepackage{multirow}
\usepackage{enumitem}
\usepackage{fancyhdr}
\usepackage{titlesec}
\usepackage{adjustbox}
\usepackage{caption}
\usepackage{draftwatermark}
\SetWatermarkText{OKernel}
\SetWatermarkScale{1.2}
\SetWatermarkColor[gray]{0.88}

% ──────────────────────────────────────────────
% Fonts
% ──────────────────────────────────────────────
\setmainfont{Helvetica Neue}
\setmonofont{Menlo}[Scale=0.85]

% ──────────────────────────────────────────────
% Colors (Dracula-inspired accents)
% ──────────────────────────────────────────────
\definecolor{accentblue}{HTML}{6272A4}
\definecolor{accentpurple}{HTML}{BD93F9}
\definecolor{accentgreen}{HTML}{50FA7B}
\definecolor{accentorange}{HTML}{FFB86C}
\definecolor{accentred}{HTML}{FF5555}
\definecolor{codebg}{HTML}{F5F5F5}
\definecolor{codeframe}{HTML}{DDDDDD}
\definecolor{sectioncolor}{HTML}{2D3748}
\definecolor{linkblue}{HTML}{3B82F6}

% ──────────────────────────────────────────────
% Hyperref
% ──────────────────────────────────────────────
\hypersetup{
    colorlinks=true,
    linkcolor=accentblue,
    urlcolor=linkblue,
    citecolor=accentblue,
    pdfauthor={OKernel Architecture Team},
    pdftitle={Aether Auto-Update System - Architecture Plan},
    pdfsubject={In-Application Auto-Update Architecture},
}

% ──────────────────────────────────────────────
% Code Listings
% ──────────────────────────────────────────────
\lstset{
    backgroundcolor=\color{codebg},
    basicstyle=\ttfamily\scriptsize,
    breaklines=true,
    breakatwhitespace=false,
    frame=single,
    rulecolor=\color{codeframe},
    framesep=4pt,
    tabsize=2,
    showstringspaces=false,
    captionpos=b,
    aboveskip=8pt,
    belowskip=8pt,
    xleftmargin=4pt,
    xrightmargin=4pt,
    keywordstyle=\color{accentpurple}\bfseries,
    stringstyle=\color{accentgreen},
    commentstyle=\color{accentblue}\itshape,
}

\lstdefinelanguage{swift}{
    keywords={struct, class, var, let, func, enum, case, static, import, return, guard, if, else, init, self, true, false, nil, @Published, @ObservedObject},
    morestring=[b]",
    morecomment=[l]{//},
    morecomment=[s]{/*}{*/},
}

\lstdefinelanguage{json}{
    morestring=[b]",
    literate=
     *{:}{{{\color{accentblue}:}}}{1}
      {,}{{{\color{accentblue},}}}{1}
      {\{}{{{\color{accentblue}\{}}}{1}
      {\}}{{{\color{accentblue}\}}}}{1}
      {[}{{{\color{accentblue}[}}}{1}
      {]}{{{\color{accentblue}]}}}{1},
}

% ──────────────────────────────────────────────
% Section Formatting
% ──────────────────────────────────────────────
\titleformat{\section}
  {\Large\bfseries\color{sectioncolor}}
  {\thesection.}{0.5em}{}[\vspace{-4pt}\textcolor{accentblue}{\rule{\columnwidth}{1pt}}]

\titleformat{\subsection}
  {\large\bfseries\color{sectioncolor}}
  {\thesubsection}{0.5em}{}

\titleformat{\subsubsection}
  {\normalsize\bfseries\color{accentblue}}
  {\thesubsubsection}{0.5em}{}

\titlespacing*{\section}{0pt}{16pt}{8pt}
\titlespacing*{\subsection}{0pt}{12pt}{4pt}
\titlespacing*{\subsubsection}{0pt}{8pt}{4pt}

% ──────────────────────────────────────────────
% Header / Footer
% ──────────────────────────────────────────────
\pagestyle{fancy}
\fancyhf{}
\fancyhead[L]{\small\textcolor{accentblue}{Aether Auto-Update Architecture}}
\fancyhead[R]{\small\textcolor{accentblue}{OKernel}}
\fancyfoot[C]{\small\thepage}
\renewcommand{\headrulewidth}{0.4pt}
\renewcommand{\headrule}{\hbox to\headwidth{\color{accentblue}\leaders\hrule height \headrulewidth\hfill}}

% ──────────────────────────────────────────────
% Custom Commands
% ──────────────────────────────────────────────
\newcommand{\fullwidthfig}[3]{%
  \begin{figure}[H]
    \centering
    \includegraphics[width=\textwidth,height=0.82\textheight,keepaspectratio]{#1}
    \caption{#2}
    \label{fig:#3}
  \end{figure}
}

\newcommand{\inlinecode}[1]{\texttt{\small\colorbox{codebg}{#1}}}

% ──────────────────────────────────────────────
% Document
% ──────────────────────────────────────────────
\begin{document}

% ──────────────────────────────────────────────
% TITLE PAGE (full width)
% ──────────────────────────────────────────────
\thispagestyle{empty}

\vspace*{2cm}

\begin{center}

{\fontsize{14}{16}\selectfont\textcolor{accentblue}{\textbf{OKERNEL ENGINEERING}}}

\vspace{0.3cm}
\textcolor{accentblue}{\rule{0.6\textwidth}{1.5pt}}

\vspace{1.5cm}

{\fontsize{28}{34}\selectfont\textbf{Aether Auto-Update System}}

\vspace{0.5cm}

{\fontsize{16}{20}\selectfont\textcolor{sectioncolor}{Architecture \& Implementation Plan}}

\vspace{2cm}

{\fontsize{12}{16}\selectfont\textcolor{accentblue}{In-Application Auto-Update for Aether Terminal Emulator\\Without Reinstallation or Gatekeeper Re-Approval}}

\vspace{3cm}

\begin{tabular}{rl}
\textbf{Status:} & Proposed \\[4pt]
\textbf{Author:} & OKernel Architecture Team \\[4pt]
\textbf{Date:} & March 19, 2026 \\[4pt]
\textbf{Version:} & 1.0 \\[4pt]
\textbf{Classification:} & Internal Engineering Document \\[4pt]
\textbf{Repository:} & \texttt{Vaiditya2207/OKernel} \\[4pt]
\textbf{Component:} & \texttt{apps/aether} \\
\end{tabular}

\vfill

\textcolor{accentblue}{\rule{0.6\textwidth}{0.5pt}}

\vspace{0.3cm}
{\small\textcolor{accentblue}{OKernel --- Cycle-Accurate Execution Visualization Platform}}\\
{\small\textcolor{accentblue}{\texttt{https://github.com/Vaiditya2207/OKernel}}}

\end{center}

\newpage

% ──────────────────────────────────────────────
% TABLE OF CONTENTS
% ──────────────────────────────────────────────
\thispagestyle{empty}
\tableofcontents
\newpage

% ──────────────────────────────────────────────
% BODY (two column)
% ──────────────────────────────────────────────
\setcounter{page}{1}

% ══════════════════════════════════════════════
\section{Executive Summary}
\label{sec:executive-summary}

The system enables Aether to check for updates, download new versions, and apply them \textbf{without reinstalling} --- avoiding macOS Gatekeeper re-approval. The update replaces the internals of the existing \texttt{.app} bundle in-place, requiring only an app restart. Users have full control via a config toggle in \inlinecode{\textasciitilde/.config/aether/config.json} or \inlinecode{config.toml}.

\subsection{Key Principles}

\begin{itemize}[nosep, leftmargin=*]
  \item \textbf{Zero friction} --- one-click update from inside the app
  \item \textbf{No Gatekeeper} --- in-place bundle patching, not a new install
  \item \textbf{User control} --- config toggle to enable/disable entirely
  \item \textbf{Resilient} --- backup before install, rollback on failure
  \item \textbf{Theme-aware} --- all UI matches the user's chosen theme
\end{itemize}


% ══════════════════════════════════════════════
\section{Why This Works Without Gatekeeper}
\label{sec:gatekeeper}

macOS Gatekeeper quarantine is applied to the \textbf{outer \texttt{.app} bundle} when first downloaded/opened. Once the user approves it, the \texttt{com.apple.quarantine} extended attribute is cleared. Replacing files \textit{inside} the bundle (executables, frameworks, resources) does \textbf{not} re-trigger Gatekeeper --- the approval lives on the \texttt{.app} directory itself.

\textbf{Critical constraint:} The bundle must be re-signed ad-hoc after patching. The \texttt{codesign} CLI is available on all macOS systems. This is the same technique used by Sparkle, VS Code, and iTerm2.

\begin{lstlisting}[title=Gatekeeper Behavior]
[YES] First install from DMG
        -> User approves once
[YES] In-place binary replace
        -> No re-approval needed
[YES] Ad-hoc re-sign after patch
        -> codesign --force --deep --sign -
[NO]  New .app from new DMG
        -> Gatekeeper triggers again
           (we avoid this)
\end{lstlisting}


% ══════════════════════════════════════════════
\section{System Architecture}
\label{sec:architecture}

The auto-update system spans three domains: the \textbf{SysCore backend} (artifact storage + API), the \textbf{Aether client} (check/download/install), and \textbf{GitHub Actions CI/CD} (build + publish).

\fullwidthfig{diagrams/architecture-light.png}{System architecture showing the three domains: Aether client, SysCore backend, and GitHub Actions CI/CD. Arrows indicate data flow for version checks, downloads, and artifact publishing.}{architecture}

\subsection{Component Responsibilities}

\begin{tabularx}{\columnwidth}{lX}
\toprule
\textbf{Component} & \textbf{Role} \\
\midrule
SysCore API & Serves version metadata and bundle tarballs \\
UpdateManager & Client-side singleton owning the full update lifecycle \\
UpdateBadgeView & Title bar indicator when update is available \\
UpdateModalView & Download/install progress overlay \\
UpdateLetterView & Post-update ``What's New'' screen \\
ConfigManager & Reads \texttt{update.enabled} and settings \\
GitHub Actions & Builds \texttt{.tar.gz} bundles and publishes to SysCore \\
\bottomrule
\end{tabularx}


% ══════════════════════════════════════════════
\section{Backend Changes (SysCore)}
\label{sec:backend}

\subsection{New: GET /api/v1/aether/latest}

Lightweight version-check endpoint. The existing \texttt{GET /api/v1/aether} returns ALL versions --- too heavy for periodic polling.

\textbf{File:} \texttt{syscore/src/server/aether.rs}

\begin{lstlisting}[language=json, title=Response]
{
  "version": "v0.4.0",
  "description": "Performance improvements",
  "changelog": "- Fixed M3 rendering\n
                 - 2x faster startup",
  "release_date": "2026-03-19T00:00:00Z",
  "size": 28456000
}
\end{lstlisting}

Reuses \texttt{list\_handlers} logic but returns only the first (latest by \texttt{release\_date}) entry. Returns \texttt{404} if no versions exist.

\textbf{Route registration} in \texttt{lib.rs}:
\begin{lstlisting}[language=C]
.route("/api/v1/aether/latest",
       get(latest_handler))
\end{lstlisting}

\subsection{New: GET /api/v1/aether/download/bundle}

Serves a \texttt{.tar.gz} of the \textbf{app bundle contents} (not a DMG). A tarball can be extracted in-place --- a DMG requires mounting and manual drag-to-Applications.

\begin{itemize}[nosep, leftmargin=*]
  \item CI/CD creates both DMG (fresh installs) and \texttt{Aether-bundle-\{ver\}.tar.gz} (updates)
  \item Tarball contains \texttt{Contents/} from inside \texttt{Aether.app/}
  \item Stored alongside DMG in \texttt{storage/aether/\{version\}/}
\end{itemize}

\textbf{Query params:} \texttt{?v=VERSION} or \texttt{?v=latest}

\textbf{Response:} Streamed \texttt{application/gzip} binary

\subsection{Upload Handler Update}

Modify \texttt{upload\_handler} to accept an optional second file field \texttt{bundle} (the \texttt{.tar.gz}). The \texttt{metadata.json} gains:

\begin{lstlisting}[language=json, title=Extended Metadata]
{
  "version": "v0.4.0",
  "filename": "Aether_Installer.dmg",
  "size": 28456000,
  "bundle_filename":
    "Aether-bundle-v0.4.0.tar.gz",
  "bundle_size": 24000000
}
\end{lstlisting}

\subsection{CI/CD Change}

\textbf{File:} \texttt{.github/workflows/aether-release.yml}

After building the \texttt{.app} bundle:
\begin{lstlisting}[language=bash]
cd build_artifact
tar -czf "Aether-bundle-${VERSION}.tar.gz" \
    -C Aether.app Contents
\end{lstlisting}

Upload both DMG and tarball in the multipart POST to SysCore.


% ══════════════════════════════════════════════
\section{Configuration}
\label{sec:config}

\subsection{Config Schema}

Add an \texttt{update} section to \texttt{AetherConfig}. Both JSON and TOML are supported.

\begin{lstlisting}[language=json, title=JSON (~/.config/aether/config.json)]
{
  "update": {
    "enabled": true,
    "check_interval": 3600,
    "server_url":
      "https://api.hackmist.tech",
    "channel": "stable"
  }
}
\end{lstlisting}

\begin{lstlisting}[title=TOML (~/.config/aether/config.toml)]
[update]
enabled = true
check_interval = 3600
server_url = "https://api.hackmist.tech"
channel = "stable"
\end{lstlisting}

\subsection{Config Fields}

\begin{tabularx}{\columnwidth}{lllX}
\toprule
\textbf{Field} & \textbf{Type} & \textbf{Default} & \textbf{Description} \\
\midrule
\texttt{enabled} & bool & \texttt{true} & Master toggle \\
\texttt{check\_interval} & int & \texttt{3600} & Seconds between checks (min: 300) \\
\texttt{server\_url} & string & hackmist & SysCore base URL \\
\texttt{channel} & string & \texttt{stable} & Future: beta/nightly \\
\bottomrule
\end{tabularx}

\subsection{Swift Config Struct}

\textbf{File:} \texttt{ConfigManager.swift}

\begin{lstlisting}[language=swift]
struct UpdateConfig: Codable {
    var enabled: Bool = true
    var checkInterval: Int = 3600
    var serverUrl: String =
      "https://api.hackmist.tech"
    var channel: String = "stable"

    enum CodingKeys: String, CodingKey {
        case enabled
        case checkInterval
                       = "check_interval"
        case serverUrl = "server_url"
        case channel
    }
}
\end{lstlisting}

Add \inlinecode{var update: UpdateConfig = UpdateConfig()} to \texttt{AetherConfig}.


% ══════════════════════════════════════════════
\section{Client-Side Architecture}
\label{sec:client}

\subsection{UpdateManager}

\textbf{File:} \texttt{Models/UpdateManager.swift}

Singleton \texttt{ObservableObject} that owns the entire update lifecycle.

\begin{lstlisting}[language=swift]
class UpdateManager: ObservableObject {
    static let shared = UpdateManager()

    @Published var state: UpdateState
                                   = .idle
    @Published var availableVersion:
                   AetherVersionInfo?
    @Published var downloadProgress:
                   Double = 0.0
    @Published var statusMessage:
                   String = ""

    enum UpdateState {
        case idle
        case checking
        case available(String)
        case downloading
        case readyToInstall
        case installing
        case failed(String)
        case restartRequired
    }

    func checkForUpdates() { }
    func downloadUpdate() { }
    func applyUpdate() { }
    func restartApp() { }
    func schedulePeriodicCheck() { }
    func cancelDownload() { }
}
\end{lstlisting}

\subsection{Version Comparison}

Semantic versioning comparison (major.minor.patch numerically). \textbf{Never string comparison.}

\begin{lstlisting}[language=swift]
struct SemanticVersion: Comparable {
    let major: Int
    let minor: Int
    let patch: Int

    init?(_ string: String) {
        let cleaned = string
            .hasPrefix("v")
            ? String(string.dropFirst())
            : string
        let parts = cleaned
            .split(separator: ".")
            .compactMap { Int($0) }
        guard parts.count == 3
            else { return nil }
        (major, minor, patch) =
            (parts[0], parts[1], parts[2])
    }

    static func < (lhs: Self,
                    rhs: Self) -> Bool {
        (lhs.major, lhs.minor, lhs.patch)
        < (rhs.major, rhs.minor, rhs.patch)
    }
}
\end{lstlisting}

\subsection{Network Resilience}

\begin{itemize}[nosep, leftmargin=*]
  \item \textbf{Timeout:} 10s for version check, 120s for download
  \item \textbf{Background failure:} Silent retry at next interval
  \item \textbf{Manual check failure:} Error in modal
  \item \textbf{Download failure:} Error + retry button in modal
\end{itemize}


% ══════════════════════════════════════════════
\section{Update State Machine}
\label{sec:state-machine}

The \texttt{UpdateManager} follows a strict state machine. Every UI element is driven by this state.

\fullwidthfig{diagrams/state-machine-light.png}{Update state machine. Transitions are triggered by timer events, user actions, and network responses. The \texttt{Disabled} state is reached when \texttt{update.enabled = false} in config.}{state-machine}

\subsection{State Transitions}

\begin{tabularx}{\columnwidth}{llX}
\toprule
\textbf{From} & \textbf{To} & \textbf{Event} \\
\midrule
idle & checking & Timer / manual \\
checking & available & Remote > Local \\
checking & idle & Remote <= Local / error \\
available & downloading & User clicks badge \\
downloading & readyToInstall & 100\% complete \\
downloading & failed & Error \\
downloading & idle & User cancels \\
readyToInstall & installing & ``Restart \& Update'' \\
readyToInstall & restartPending & ``Later'' \\
installing & restartRequired & Success \\
installing & failed & Error \\
failed & checking & Retry \\
failed & idle & Dismiss \\
\bottomrule
\end{tabularx}


% ══════════════════════════════════════════════
\section{Complete Update Flow}
\label{sec:update-flow}

End-to-end flow from app launch to post-update experience:

\begin{figure}[H]
  \centering
  \includegraphics[height=0.88\textheight,keepaspectratio]{diagrams/update-flow-light.png}
  \caption{Complete update flow from app launch through version check, download, in-place install, and post-update ``What's New'' experience.}
  \label{fig:update-flow}
\end{figure}


% ══════════════════════════════════════════════
\section{Install Sequence}
\label{sec:install-sequence}

Detailed interaction between all components during the update process:

\fullwidthfig{diagrams/install-sequence-light.png}{Sequence diagram showing interactions between User, UpdateBadgeView, UpdateModalView, UpdateManager, SysCore API, Filesystem, and Aether.app during the update lifecycle.}{install-sequence}

\subsection{Critical Path Steps}

\begin{enumerate}[nosep, leftmargin=*]
  \item Verify download integrity (file size matches \texttt{metadata.size})
  \item Locate current app bundle: \texttt{Bundle.main.bundlePath}
  \item \textbf{Backup:} Copy \texttt{Contents/} to \texttt{\textasciitilde/Library/Application Support/Aether/backup/}
  \item \textbf{Extract:} \texttt{tar -xzf \{download\} -C \{bundle\}/}
  \item \textbf{Re-sign:} \texttt{codesign --force --deep --sign - \{bundle\}}
  \item Set \texttt{pendingUpdateVersion} in UserDefaults
  \item Clean up download artifacts
  \item \textbf{Restart:} \texttt{NSWorkspace.open(bundleURL)} + \texttt{NSApp.terminate()}
\end{enumerate}

\textbf{Why \texttt{tar} via Process?} Swift's built-in archive support is limited. \texttt{/usr/bin/tar} is on every macOS, handles \texttt{.tar.gz} natively. Run via \texttt{Process()}.

\textbf{Why ad-hoc re-sign?} Replacing the executable invalidates the code signature. Ad-hoc signing creates a local signature without a Developer ID. This is what the current \texttt{build\_release.sh} already does.

\textbf{Rollback:} If step 4 or 5 fails, restore from backup immediately.


% ══════════════════════════════════════════════
\section{UI Components}
\label{sec:ui}

\subsection{Title Bar Layout}

The update badge sits on the right side of the existing 28pt custom title bar:

\fullwidthfig{diagrams/titlebar-layout-light.png}{Title bar layout showing traffic lights (left), centered session title, and the update badge pill button (right). The tab bar sits below.}{titlebar}

\subsection{Badge Specification}

\begin{tabularx}{\columnwidth}{lX}
\toprule
\textbf{Property} & \textbf{Value} \\
\midrule
Shape & Capsule (fully rounded pill) \\
Font & System 10pt, medium weight, monospaced \\
Padding & Horizontal 8pt, vertical 3pt \\
Background & Theme \texttt{selection} at 60\% opacity \\
Text color & Theme \texttt{foreground} color \\
Hover & Background opacity increases to 90\% \\
Animation & 0.3s fade-in \\
Content & ``Update v\{X.Y.Z\}'' + SF Symbol \\
Visibility & Hidden when disabled or up-to-date \\
\bottomrule
\end{tabularx}

\textbf{Integration point} in \texttt{AetherApp.swift} (lines 48--51):

\begin{lstlisting}[language=swift]
HStack {
    Spacer()
    if configManager.config.update.enabled {
        UpdateBadgeView(
            updateManager: updateManager,
            theme: configManager.config.colors
        )
        .padding(.trailing, 8)
    }
}
.frame(height: 28)
\end{lstlisting}

\subsection{Update Modal}

Centered overlay with glass morphism styling.

\fullwidthfig{diagrams/modal-states-light.png}{Modal states flow: Checking (spinner) leads to Downloading (progress bar), which leads to Ready to Install (restart/later buttons) or Failed (retry/dismiss). Installing shows a non-dismissable spinner.}{modal-states}

\begin{tabularx}{\columnwidth}{lX}
\toprule
\textbf{Property} & \textbf{Value} \\
\midrule
Background & \texttt{VisualEffectView(.hudWindow)} + dark overlay \\
Corner radius & 12pt \\
Width & 320pt fixed \\
Height & Intrinsic (content-driven) \\
Buttons & Capsule, theme-aware colors \\
Progress bar & Capsule, accent fill, secondary track \\
Animation & Spring entrance, opacity exit \\
Dismiss & Click outside (not during install) \\
\bottomrule
\end{tabularx}

\subsection{State-Driven Modal Content}

\begin{tabularx}{\columnwidth}{lX}
\toprule
\textbf{State} & \textbf{Content} \\
\midrule
\texttt{.checking} & Spinner + ``Checking...'' \\
\texttt{.downloading} & Progress bar + size + Cancel \\
\texttt{.readyToInstall} & ``Restart Now'' + ``Later'' \\
\texttt{.installing} & Spinner + ``Do not quit'' \\
\texttt{.failed} & Error + Retry + Dismiss \\
\bottomrule
\end{tabularx}

\subsection{Post-Update ``What's New'' View}

Shown after restart from an update. Reuses \texttt{StartupView}'s visual language.

\begin{enumerate}[nosep, leftmargin=*]
  \item \textbf{Stage 1:} OKernel logo (2.5s, same as first-run)
  \item \textbf{Stage 2:} ``What's New in Aether \{version\}'' --- changelog with glass styling
  \item \textbf{[CONTINUE]} button leads to normal terminal
\end{enumerate}

\textbf{Trigger:} Check \texttt{pendingUpdateVersion} in UserDefaults on launch. If present, show this flow. Clear after display.


% ══════════════════════════════════════════════
\section{File System Layout}
\label{sec:filesystem}

\fullwidthfig{diagrams/file-tree-light.png}{File system layout across three locations: user config (\textasciitilde/.config/aether), application support (downloads and backups), and the Aether.app bundle that gets patched in-place.}{filesystem}

\begin{lstlisting}[title=Directory Structure]
~/Library/Application Support/Aether/
  sessions.plist        (existing)
  updates/
    v0.4.0/
      Aether-bundle-v0.4.0.tar.gz
  backup/
    Contents/           (pre-update copy)

~/.config/aether/
  config.json           (gains "update")
  config.toml

/Applications/Aether.app/
  Contents/
    MacOS/AetherApp     (replaced)
    Frameworks/         (replaced)
    Resources/          (replaced)
    Info.plist          (updated)
\end{lstlisting}


% ══════════════════════════════════════════════
\section{Security Considerations}
\label{sec:security}

\begin{tabularx}{\columnwidth}{lX}
\toprule
\textbf{Concern} & \textbf{Mitigation} \\
\midrule
MITM & HTTPS via \texttt{api.hackmist.tech}. ATS enforced. \\
Tampered binary & Verify size vs \texttt{metadata.size}. Phase 2: SHA-256 checksum. \\
Privilege escalation & No \texttt{sudo}. Writes only to own bundle + user dirs. \\
Downgrade attack & Only update if remote > local (semver). \\
Interrupted install & Backup exists. Restore on failure. \\
Disk space & Check before download (need 2x update size). \\
Read-only bundle & Detect. Show: ``Move Aether to writable location.'' \\
Running from DMG & Detect read-only FS. Show: ``Install Aether first.'' \\
\bottomrule
\end{tabularx}


% ══════════════════════════════════════════════
\section{State Persistence}
\label{sec:persistence}

\begin{tabularx}{\columnwidth}{llX}
\toprule
\textbf{Key} & \textbf{Type} & \textbf{Purpose} \\
\midrule
\texttt{pendingUpdateVersion} & String? & Set before restart \\
\texttt{pendingUpdateChangelog} & String? & Changelog for ``What's New'' \\
\texttt{lastUpdateCheck} & Date? & Last successful check \\
\texttt{updateDownloadedVersion} & String? & Completed download \\
\bottomrule
\end{tabularx}


% ══════════════════════════════════════════════
\section{Implementation Phases}
\label{sec:phases}

\fullwidthfig{diagrams/phase-gantt-light.png}{Gantt chart of 7 implementation phases: Foundation, Update Check, Title Bar Badge, Download + Modal, Install + Restart, Post-Update UX, and Polish + Tests.}{gantt}

\subsection{Phase 1: Foundation}

\begin{enumerate}[nosep, leftmargin=*]
  \item Add \texttt{GET /api/v1/aether/latest} to \texttt{aether.rs}
  \item Add \texttt{UpdateConfig} struct to \texttt{ConfigManager.swift}
  \item Update \texttt{config.example.json} and \texttt{.toml}
  \item Update config decoder with defaults
\end{enumerate}

\subsection{Phase 2: Update Check}

\begin{enumerate}[nosep, leftmargin=*, start=5]
  \item Create \texttt{UpdateManager.swift} singleton
  \item Implement \texttt{checkForUpdates()} with URLSession
  \item Implement semver comparison
  \item Implement \texttt{schedulePeriodicCheck()} with Timer
  \item Wire into \texttt{AetherApp.swift} on launch
\end{enumerate}

\subsection{Phase 3: Title Bar Badge}

\begin{enumerate}[nosep, leftmargin=*, start=10]
  \item Create \texttt{UpdateBadgeView.swift}
  \item Integrate into title bar ZStack
  \item Bind to \texttt{UpdateManager.state}
\end{enumerate}

\subsection{Phase 4: Download + Modal}

\begin{enumerate}[nosep, leftmargin=*, start=13]
  \item Create \texttt{UpdateModalView.swift}
  \item Implement download with progress delegate
  \item Implement \texttt{cancelDownload()}
  \item Wire modal to badge tap + state
\end{enumerate}

\subsection{Phase 5: Install + Restart}

\begin{enumerate}[nosep, leftmargin=*, start=17]
  \item Implement \texttt{applyUpdate()} (backup, extract, re-sign)
  \item Implement \texttt{restartApp()} (NSWorkspace + terminate)
  \item Add tarball generation to CI/CD
  \item Add bundle download endpoint
\end{enumerate}

\subsection{Phase 6: Post-Update Experience}

\begin{enumerate}[nosep, leftmargin=*, start=21]
  \item Create \texttt{UpdateLetterView.swift}
  \item Check \texttt{pendingUpdateVersion} on launch
  \item Show logo then changelog sequence
\end{enumerate}

\subsection{Phase 7: Polish + Tests}

\begin{enumerate}[nosep, leftmargin=*, start=24]
  \item Disk space check before download
  \item Backup restoration on failure
  \item SHA-256 checksum verification
  \item Auto-cleanup of old downloads/backups
  \item Write all unit, integration, and E2E tests
\end{enumerate}


% ══════════════════════════════════════════════
\section{Test Plan}
\label{sec:tests}

\subsection{Unit Tests}

\textbf{File:} \texttt{Tests/AetherAppTests/UpdateManagerTests.swift}

\begin{tabularx}{\columnwidth}{lX}
\toprule
\textbf{ID} & \textbf{Verifies} \\
\midrule
U1 & \texttt{v1.0.0 > v0.9.9} (major) \\
U2 & \texttt{v0.4.0 > v0.3.0} (minor) \\
U3 & \texttt{v0.3.1 > v0.3.0} (patch) \\
U4 & \texttt{v0.3.0 == v0.3.0} (no update) \\
U5 & \texttt{v0.2.0 < v0.3.0} (no downgrade) \\
U6 & Handles \texttt{v} prefix and bare versions \\
U7 & Default \texttt{UpdateConfig.enabled = true} \\
U8 & Disabled config prevents checks \\
U9 & Interval < 300 clamped to 300 \\
U10 & Missing \texttt{update} key uses defaults \\
U11 & TOML \texttt{[update]} parses correctly \\
U12 & Full state machine transitions \\
U13 & Progress calculation correctness \\
U14 & Backup path construction \\
U15 & Download path construction \\
\bottomrule
\end{tabularx}

\vspace{6pt}
\textbf{File:} \texttt{syscore/tests/api\_tests.rs}

\begin{tabularx}{\columnwidth}{lX}
\toprule
\textbf{ID} & \textbf{Verifies} \\
\midrule
U16 & \texttt{/latest} returns newest by date \\
U17 & \texttt{/latest} returns 404 when empty \\
U18 & Response has all required fields \\
U19 & Bundle download serves gzip \\
U20 & \texttt{?v=latest} resolves correctly \\
\bottomrule
\end{tabularx}

\subsection{Integration Tests}

\textbf{File:} \texttt{Tests/AetherAppTests/UpdateIntegrationTests.swift}

\begin{tabularx}{\columnwidth}{lX}
\toprule
\textbf{ID} & \textbf{Verifies} \\
\midrule
I1 & Mock server check transitions to \texttt{.available} \\
I2 & Same version stays \texttt{.idle} \\
I3 & Unreachable server: no crash, retries \\
I4 & Malformed JSON: graceful handling \\
I5 & Download lands at correct path + size \\
I6 & Cancel cleans up partial file \\
I7 & Progress callbacks fire 0.0 to 1.0 \\
I8 & Runtime toggle invalidates timer \\
I9 & Pending update persisted in UserDefaults \\
I10 & Backup + corrupt + restore works \\
\bottomrule
\end{tabularx}

\subsection{End-to-End Tests}

\textbf{File:} \texttt{Tests/AetherAppTests/UpdateE2ETests.swift}

\textbf{Prerequisites:} Local SysCore on \texttt{127.0.0.1:3001}, two test \texttt{.tar.gz} bundles (v0.3.0 and v0.4.0).

\subsubsection{test\_full\_update\_lifecycle}

\begin{enumerate}[nosep, leftmargin=*]
  \item \textbf{Setup:} Upload v0.3.0 and v0.4.0 to SysCore. Create temp \texttt{.app} bundle at v0.3.0.
  \item \textbf{Check:} \texttt{checkForUpdates()} transitions to \texttt{.available("v0.4.0")}
  \item \textbf{Download:} \texttt{downloadUpdate()} completes; file exists with correct size
  \item \textbf{Install:} \texttt{applyUpdate()} updates Info.plist to v0.4.0; backup exists
  \item \textbf{Verify:} \texttt{pendingUpdateVersion} set; ``What's New'' would trigger
  \item \textbf{Cleanup:} Delete temp dirs and test versions
\end{enumerate}

\subsubsection{test\_update\_with\_disabled\_config}

\begin{enumerate}[nosep, leftmargin=*]
  \item Set \texttt{update.enabled = false}
  \item Call \texttt{checkForUpdates()}
  \item Assert: zero network requests (mock URLProtocol)
  \item Assert: state == \texttt{.idle}, badge not visible
\end{enumerate}

\subsubsection{test\_update\_rollback\_on\_failure}

\begin{enumerate}[nosep, leftmargin=*]
  \item Create temp bundle (v0.3.0), download v0.4.0
  \item Make bundle path read-only (simulate failure)
  \item \texttt{applyUpdate()} results in \texttt{.failed}
  \item Backup restored --- bundle still v0.3.0
\end{enumerate}

\subsubsection{test\_interrupted\_download\_cleanup}

\begin{enumerate}[nosep, leftmargin=*]
  \item Start download, cancel at ~50\%
  \item Partial file cleaned up
  \item Restart download from scratch succeeds
\end{enumerate}


% ══════════════════════════════════════════════
\section{Risks and Mitigations}
\label{sec:risks}

\begin{tabularx}{\columnwidth}{lllX}
\toprule
\textbf{Risk} & \textbf{P} & \textbf{I} & \textbf{Mitigation} \\
\midrule
Read-only bundle path & M & H & Check write permission. Show error. \\
Running from DMG & L & H & Detect RO filesystem. \\
Partial extraction & L & H & Staging dir + atomic swap + backup. \\
Misconfigured URL & L & M & Validate on load. Log. \\
Slow connection & M & M & Show speed/ETA. Cancel/retry. \\
Concurrent updates & L & M & File lock in updates dir. \\
\bottomrule
\end{tabularx}

\small{P = Probability (L/M/H), I = Impact (L/M/H)}


% ══════════════════════════════════════════════
\section{File Change Summary}
\label{sec:files}

\begin{tabularx}{\columnwidth}{llX}
\toprule
\textbf{File} & \textbf{Act} & \textbf{Description} \\
\midrule
\texttt{aether.rs} & MOD & Add \texttt{latest\_handler}, bundle download \\
\texttt{lib.rs} & MOD & Register new routes \\
\texttt{ConfigManager.swift} & MOD & Add \texttt{UpdateConfig} \\
\texttt{UpdateManager.swift} & NEW & Core update singleton \\
\texttt{UpdateBadgeView.swift} & NEW & Title bar pill button \\
\texttt{UpdateModalView.swift} & NEW & Download/install modal \\
\texttt{UpdateLetterView.swift} & NEW & ``What's New'' screen \\
\texttt{AetherApp.swift} & MOD & Wire manager, badge, modal \\
\texttt{StartupView.swift} & MOD & Add update letter stage \\
\texttt{config.example.json} & MOD & Add \texttt{update} section \\
\texttt{config.example.toml} & MOD & Add \texttt{[update]} section \\
\texttt{aether-release.yml} & MOD & Tarball generation \\
\texttt{UpdateManagerTests} & NEW & Unit tests (U1--U15) \\
\texttt{UpdateIntegrationTests} & NEW & Integration (I1--I10) \\
\texttt{UpdateE2ETests} & NEW & End-to-end tests \\
\texttt{api\_tests.rs} & MOD & Backend tests (U16--U20) \\
\bottomrule
\end{tabularx}


\end{document}
